\chapter{无穷小演算.}

1604 年，正当科学生涯的顶点，伽利略自信已经证明：在速度与所经距离成正比增长的直线运动中，运动定律恰好就是（\(x = ct^{2}\)）他在落体中发现的那条定律（{[}122 b{]}，第 X 卷，第 115–116 页）。从 1695 年到 1700 年，在莱比锡按月出版的《学术纪事》（Acta Eruditorum）中，没有一卷不刊登莱布尼茨、伯努利兄弟、洛必达侯爵的论文，它们以我们至今仍在使用的记号，处理微分学、积分学和变分法中最为各式各样的问题。因此，无穷小演算恰好是在差不多整整一个世纪的间隔之内锻造完成的，或者，如英国人后来所说，是处于最佳状态的 Calculus（“calculus”）；而近三个世纪的持续使用，也还没有完全磨钝这件无与伦比的工具。

希腊人既不曾拥有过、也不曾想象过任何与之类似的东西。如果说他们无疑知道一种代数演算，即巴比伦人的那种（他们的几何学有一部分或许只是它的转录），那也只是为了拒绝使用它；他们或许是最为辉煌的数学创造，即处理那些在我们看来会与积分学相联系的问题的方法，则严格地登记在几何学发明的领域之内。欧多克索斯在处理圆锥和棱锥的体积时给出了这一方法最初的典范，欧几里得又或多或少忠实地把它们传给了我们（{[}107{]}，第 XII 卷，命题 7 及 10）。但尤其重要的是，阿基米德的大部分工作正是献身于这些问题{[}5 b and c{]}；而且，出于一种罕见的幸运，我们今天仍能以原文读到他的大部分著作，即他以如此细心的态度用以写作的那铿锵的多利安方言原文，其中甚至包括最近重新发现的那一部，他在其中陈述了引导他得到某些最优美结果的“启发式”程序（{[}5 b{]}，第 II 卷，第 425–507 页）。因为正是在这里有着欧多克索斯“穷竭法”的弱点之一：作为一种证明方法（在假定了若干公设之后）无可指摘，它却不是一种发现的方法；它的应用必然以预先知道所要证明的结果为前提；所以，阿基米德说，“在欧多克索斯第一个给出证明的那些结果中，关于圆锥和棱锥的\ldots，不小的一部分属于德谟克利特，他是第一个陈述它们而未加证明的人”（同上出处，第 430 页）。这种情况使得细致地分析阿基米德的工作格外困难，而这种分析，说实在的，似乎还不曾有任何一位现代史学家着手进行过；因为正是由于这个缘故，我们不知道他对联结他所处理的各个问题的那些关系线意识到何种程度（这些联系，我们将用“同一个积分以多种不同的几何面貌反复出现”的说法来表达），也不知道他可能赋予它们怎样的重要性。例如，让我们考虑下面这些问题，第一个由欧多克索斯解决，其余的由阿基米德解决：棱锥的体积、抛物线弓形的面积、三角形的重心，以及所谓阿基米德螺线的面积（以极坐标表示即 \(\rho = c\omega\)）；它们全都依赖于积分 \(\int x^2 dx\)，而且，丝毫不背离穷竭法的精神，它们全都可以化归为形如 \(\sum_n an^2\) 的“黎曼和”的计算。事实上，阿基米德正是以这种方式处理螺线的（{[}5 b{]}，第 II 卷，第 1–121 页），他借助一条引理，把问题归结为写出

\[
N ^ {3} <   3 \sum_ {n = 1} ^ {N} n ^ {2} = N ^ {3} + N ^ {2} + \sum_ {n = 1} ^ {N} n <   (N + 1) ^ {3}.
\]

至于三角形的重心，他证明（用穷竭法，借助分解成平行的细条）它位于每条中线上，因而位于它们的交点（{[}5 b{]}，第 II 卷，第 261–315 页）。对于抛物线，他给出三种程序：第一种是启发式的，只图“给结果以某种可信性”，它借助一个静力学的论证把问题化归为三角形的重心，在论证过程中他不迟疑把抛物线弓形看作平行于轴的无穷多个线段之和（{[}5 b{]}，第 II 卷，第 435–439 页）；另一种方法依据一个类似的原则，但以穷竭法完全严格地加以陈述（{[}5 b{]}，第 II 卷，第 261–315 页）；最后一种证明极为巧妙但适用范围较小，它借助抛物线的特殊性质把所求面积给出为几何级数之和。没有任何迹象表明这些问题与棱锥的体积之间有联系；甚至有明文规定（{[}5 b{]}，第 II 卷，第 8 页）：与螺线有关的那些问题，同他在同一篇引言中恰好谈到的另一些关于球和旋转抛物面的问题“毫无共同之处”，而在后一类问题中恰有一个（旋转抛物面的体积）化归为积分 \(\int x dx\)。

从这些例子可以看出，撇开特殊技巧的使用不谈，穷竭法的原则如下：通过分解成“黎曼和”，得到所考虑的量的上界与下界，再把这些界值与为该量给出的表达式直接比较，或者与一个已经解决的类似问题的相应界值比较。这种比较（由于不可能使用负数而必然分两步进行）以套语引入：“若不然，实际上它就会或者较大或者较小；假若可能，设它较大，等等；假若可能，设它较小，等等。”由此得出了 17 世纪的学者们给予该方法的名称：“归谬的”或“通过化归于荒谬”（“ἀπαγωγὴ εἰς ἀδύνατον”）。阿基米德正是以一种类似的形式作出螺线的切线的确定（{[}5 b{]}，第 II 卷，第 62–76 页），这是一个孤立的结果，也是我们能够引作“微分学”古代渊源的唯一结果——除了圆锥曲线的切线的相对容易的确定以及若干极大极小问题之外。事实上，就“积分”而言，一个广阔的研究领域摆在了希腊数学家面前，不仅有面积和体积的理论，还有静力学和流体静力学；而由于缺乏动力学，他们几乎没有机会认真地着手处理微分。诚然，阿基米德给出了他的螺线的一种动力学定义；而且，既然我们不知道他怎样会被引向对螺线切线的认识，我们就有权自问，他是否对运动的合成有某种想法。但若是那样，他难道不会把这样一个有力的方法应用到其他同类问题上去吗？更可能的是，他必定使用了某种过渡到极限的启发式程序，这是他所知道的关于圆锥曲线的结果本可以提示的；这些结果当然具有本质上更为简单的性质，因为人们可以作出直线与圆锥曲线的交点，从而确定这些点重合的条件。至于切线的定义，它被设想为一条在曲线上某点的邻域内把曲线完全留在同一侧的直线；它的存在是被假定的，而且还假定每条曲线都由凸弧组成；在这些条件下，要证明一条直线与一条曲线相切，必须证明某些不等式，而这当然是以最完善的准确性来完成的。

从严格性的观点看，阿基米德的方法无可指摘；而且在 17 世纪，当最一丝不苟的数学家们想把一个被认为特别微妙的结果置于一切怀疑之外时，他们给出的仍是一种“归谬的”证明（{[}109{]}，第 I 卷，第 211–254 页；法文译本，第 III 卷，第 181–215 页，及{[}244{]}，第 VIII 卷，第 249–282 页）。至于其丰饶多产，阿基米德的工作本身便是充分的见证。但是，要想有权在那里看到一种“积分学”，就必须在几何面貌的多样性之中，显示出按作为基础的“积分”的本性对问题作分类的某种轮廓。在 17 世纪，正如我们将要看到的，对这样一种分类的寻求渐渐成为几何学家们的主要关切之一；如果在阿基米德那里找不到这种分类的任何痕迹，这难道不是一个标志，表明这样的思辨在他看来会是过分“抽象”的，表明他反而自愿地、在每个场合都尽可能贴近他所研究的图形的特殊性质吗？而我们难道不该得出结论说，这件令人赞叹的工作——创建者们也承认，积分学整个是从它那里取得的——在某种意义上是积分学的反面吗？

况且，在数学中，在发现与证明之间让人挖一道鸿沟，是不会不受惩罚的。在顺利的时代，数学家只须几乎像他构思时那样写下他的思想而并不放弃严格性；有时他甚至可以希望做到这一点，所付的代价只是语言和所采用记号的一次幸运的改变。但他常常必须听命于在两者之间作选择：不正确的、却可能多产的表述方法，和正确的、却除了歪曲他的思想之外、除非付出令人生厌的努力便不再容他表达思想的方法。两条道路无一免除危险。希腊人走了第二条，而他们的数学在其最辉煌的绽放之后几乎立即出现令人惊异的停滞，其原因或许正应在这里寻找，甚至更应在这里、而不是在罗马征服的使之贫瘠化的影响之中寻找。有人提出过一种并非显得不可能的看法：阿基米德和阿波罗尼奥斯的后继者们的口头讲授，本可以包含许多新结果，只是他们不曾认为有必要给自己加上符合通行典范的出版所要求的那种非常的努力。无论如何，正是这样的顾忌不再阻止 17 世纪的数学家们：当成批的新问题被陈述出来时，他们在对阿基米德著作的刻苦研读中寻求绕过它的手段。

虽然希腊文学和哲学的古典著作都已在意大利、由阿尔杜斯·马努提乌斯及其后继者们刊印，而且几乎全部在 1520 年以前，但直到 1544 年，才在巴塞尔、由埃尔瓦吉乌斯出版了最早的阿基米德版本，希腊文与拉丁文合刊{[}5 a{]}，此前不曾有任何拉丁文出版物为其开路；而且，与其说当时的数学家们（他们正埋头于代数研究）立即感受到了它的影响，不如说我们必须等到伽利略和开普勒——两人与其说是数学家不如说是天文学家和物理学家——这影响才变得明显起来。从这一时刻起，直到约 1670 年从未间断，在无穷小演算奠基者们的著作中，没有哪个名字比阿基米德的名字出现得更频繁。许多人翻译并注释他；从费马到巴罗，所有人都竞相引用他；所有人都宣称在那里同时找到一个典范和一种灵感的源泉。

诚然，正如我们将要看到的，这些宣言不应全部按字面理解；这正是我们在设法对这些著作作出真实解释时碰到的困难之一。史学家还必须考虑到当时科学界的组织状况：17 世纪初它仍非常不完善，而到同一世纪末，多亏学术团体和科学期刊的创立，多亏大学的巩固和发展，它终于变得与我们今天所知的非常相像。直到 1655 年都没有任何期刊，数学家们要让其工作为人所知，除通信和印书之外别无选择，而印书多半自费，或者，若能找到一位庇护人，则由庇护人出资。能承担这类工作的编辑和印刷商很稀少，而且往往相当靠不住。在这种出版所意味着的长期耽搁和无尽麻烦之后，作者还常常不得不硬着头皮面对无休止的论战，这些论战由并不总是怀着善意的对手挑起，有时以出人意料尖刻的语气进行：因为，在甚至连无穷小演算的原理都还处于普遍不明朗的境地时，每个人都不难在对手的论证中找到薄弱之处，或至少是含混而可争议之处。可以理解，在这些条件下，许多爱好平静的学者满足于把他们的方法和结果只传达给少数挑选过的朋友。有些人，尤其是某些科学的业余爱好者，例如巴黎的梅森、以及后来伦敦的柯林斯，与各国保持着大量的通信，他们在通信中东一处西一处地传播摘录，并不回避掺入他们自家酿造的蠢话。掌握着“方法”的数学家们，由于缺乏概念和一般的定义，不能把方法写成定理的形式，甚至不能以某种准确性加以表述，于是只好借助于大量特定的情形来作尝试，并且相信，要衡量方法的能力，再没有比向同行们发出挑战更好的办法了，有时还伴随着把他们的结果用密码发表。勤学的青年人四处游学，也许比今天游学得更多；一位学者的思想有时借助其学生的游历传播得比借助他本人的著作更好，但这并非不因此又成为误解的一个根源。最后，同样这些问题必然出现在大量数学家面前，其中有许多非常杰出的人物，而他们彼此对他人结果只有不完善的了解，于是优先权的主张不可能不接连不断地出现，而且加上剽窃指控的情形也不罕见。

因此，史学家必须在当时学者们的书信和私人文稿中、与其正式出版物同样多地、甚至更多地寻找他的文献。但是，例如惠更斯的书信保存了下来，并已成为典范性出版的对象{[}169 b{]}，而莱布尼茨的书信至今只以片段的方式发表，许多别的则已无可挽回地散失。至少，最近以手稿分析为基础的研究，以一种似乎无可辩驳的方式，澄清了一个曾被派性论争多少掩盖了的要点：那就是，每当那个时代的一位大数学家叙述他自己的工作、他思想的发展、他所受到的和所未受到的影响时，他都是以诚实而真诚的方式、以完全的善意行事的；\(^{1}\)这些珍贵的证言，我们拥有相当多，因而可以放心地加以使用，史学家在它们的情形下并不需要把自己变成一位预审法官。至于其余的，大多数已出现的优先权问题都完全没有意义。诚然，莱布尼茨在采用 dx 这个记号表示“微分”时，并不知道牛顿以 \(\dot{x}\) 表示“流数”已有大约十年之久；但即使他知道，又有什么要紧呢？取一个更有教益的例子：定理 \(\log x = \int \frac{dx}{x}\) 的作者是谁，日期又是什么？这个公式，正如我们刚才所写下的那样，出自莱布尼茨，因为它的两个成员都是用他的记号写出的。莱布尼茨本人和沃利斯都把它归之于格雷古瓦·德·圣万桑。后者在他的《几何著作》（Opus Geometricum）{[}135{]}（出版于 1647 年，但据他说写成于很久以前）中，只证明了下述命题的等价物：若 \(f(a,b)\) 表示双曲线弓形 \(a \leq x \leq b, 0 \leq y \leq A/x\) 的面积，则关系 \(b'/a' = (b/a)^n\) 蕴含 \(f(a'/b') = n.f(a,b)\)；对此，他的学生和注释者萨拉萨几乎立即补充{[}269{]}说，面积 \(f(a,b)\) 因此可以“取对数之位”。如果他只说了这些，而格雷古瓦本人又未对此多置一词，这难道不是因为，对当时的大多数数学家来说，对数只是些“计算辅助工具”，在数学中没有资格被引用吗？诚然，托里拆利在 1644 年的一封信{[}208{]}中谈到他对一条我们应写作 \(y = ae^{-cx}, x \geq 0\) 的曲线的研究，同时补充说，在纳皮尔（他此外还对他备加赞扬）“只追求算术用途”之处，他本人“从中引出了一种关于几何学的思辨”；而且他留下一篇关于这条曲线的、显然已为发表而准备的手稿，但它直到 1900 年仍未出版（{[}310{]}，第 I 卷，第 335–347 页）。笛卡尔还于 1639 年在“德博纳问题”的联系中遇到过同一条曲线，并描述了它而没有谈到对数（{[}85 a{]}，第 II 卷，第 514–517 页）。无论如何，J. 格雷戈里在 1667 年给出了一条用（十进）对数计算双曲线弓形面积的规则，没有引用任何人（{[}136 a{]}，转载于{[}169 a{]}，第 407–462 页）：这同时蕴含着对双曲线求积与对数之间联系的理论认识，和对“自然”对数与“十进”对数之间联系的数值认识。惠更斯的主张是否只针对最后这一点？他立即质疑了格雷戈里这一结果的新颖性（{[}169 b{]}，第 VI 卷，第 228–230 页）。这件事对我们和对他的同时代人一样不清楚；这些同时代人无论如何都清楚地感到，对数与双曲线求积之间联系的存在是一件久已知晓的事情，只是他们在这点上除了书信中的暗示或格雷古瓦·德·圣万

桑。1668 年，当布龙克尔给出（附有一个打磨过的收敛性证明，即与一个几何级数的比较）log 2 和 log (5/4) 的级数（{[}38{]}：转载于{[}218{]}，第 I 卷，第 213–218 页）时，他把它们作为双曲线相应弓形的表达式来呈现，并补充说，他所得到的数值与 2 和 5/4 的对数“成相同的比”。但在同一年，随着墨卡托（{[}220{]}，转载于{[}218{]}，第 I 卷，第 167–196 页），或更确切地说，随着沃利斯立即对墨卡托工作所作的阐述（{[}326 b{]}，转载于{[}218{]}，第 I 卷，第 219–226 页），语言改变了：既然双曲线的弓形与对数成比例，而众所周知，对数由其特征性质只确定到一个常数因子，那么没有任何东西妨碍把双曲线的弓形看作对数，其资格是“自然的”（与“人为的”或“十进的”对数相对），或者说双曲线的；这最后一步一旦跨出（墨卡托给出的 log(1+x) 的级数对此有所贡献），定理 log x = ∫ dx/x 便得到了，至多差个记号，或者毋宁说它变成了一个定义。能得出什么结论呢？只能得出：这一发现是通过几乎察觉不到的过渡完成的，而围绕此事的优先权之争，与小提琴和长号为一部交响曲中某一动机出现的准确时刻而起的争执极为相似。而且，说实在的，与此同时代的别的数学创造，费马的算术、牛顿的动力学，都带有强烈的个人印记，而 17 世纪无穷小演算的发展使我们想起的，却很是一部交响曲那种渐进而必然的发展，在那里，“时代精神”（Zeitgeist）既作作曲家又当指挥，执掌着指挥棒：每个人以特有的音色演奏自己的声部，但没有一个人是他为听众所创造的那些主题的主人，这些主题被一种精巧的对位法几乎无法解脱地缠绕在一起。因此，这一历史应当以主题分析的形式来写；我们在这里只满足于一个概略的草图，它并不奢求一丝不苟的准确。² 下面就是一番肤浅的考察所呈现的主要主题：

\begin{enumerate}
\def\labelenumi{\Alph{enumi})}
\item
  数学严格性的主题，与无穷小、不可分量或微分的主题相对照。我们已经看到，二者在阿基米德那里都占有重要的位置：第一个贯穿于他的全部工作，第二个则只在他论《方法》的著作中，而 17 世纪并不知道这部著作，因此，倘若它是被传下来而不是被重新发明的，那就只能是通过哲学的传统。无穷小的原则此外还以两种不同的形式出现，视所涉及的是“微分”还是“积分”而定。就后者而言，设先要计算的是一个平面面积：人们将用无穷多条等距的平行线把它分成无穷多条无穷小的平行细条；而每一条这样的细条都是一个矩形（尽管用相距有限距离的两条平行线所得到的有限细条，没有一条会是矩形）。同样，一个旋转体将被垂直于轴的平面分解成无穷多个具有同一无穷小高度的圆柱；\(^{3}\) 当涉及用共点的直线把一个面积分解成三角形，或者把曲线一段弧的长当作一个具有无穷多条边的多边形来讨论时，也可以使用类似的说法，等等。可以肯定，那些完全掌握运用阿基米德方法之能力的罕见数学家，如费马、帕斯卡、惠更斯、巴罗，本可以在每个特定的情形中毫无困难地用严格的证明来代替这种语言的使用；因此他们也常常指出，这种语言只是为了简洁才存在。“本来很容易”，费马说，“按阿基米德的方式给出证明\ldots{}；话既说到，只需一次性预先提醒，便可避免不断的重复\ldots{}”（{[}109{]}，第 I 卷，第 257 页）；帕斯卡同样说：“这两个方法之一与另一个只是陈述的方式不同而已”（{[}244{]}，第 VIII 卷，第 352 页）；\(^{4}\) 而已罗则带着他嘲弄式的简洁说：“longior discursus apagogicus adhiberi possit, sed quorsum?”（本可以用一段更长的归谬论证来处理，但何苦呢？）（{[}16 b{]}，第 251 页）。费马甚至看来非常谨慎，不肯拿出任何他不能如此证成的东西，并因此自我限定：除暗示外、或以“方法”的形式，不宣布任何一般的结果；巴罗虽然同样谨慎，却没有那么多的顾忌。至于他们同时代人的大多数，至少可以说，严格性不是他们的主要关切，而阿基米德的名字多半只是一种掩饰，用来遮盖无疑颇有价值的货物，只是阿基米德肯定不会愿意为它承担责任。在涉及微分的场合就更是如此。如果说，当涉及曲线的求长时，曲线被当作一个具有无穷多条边的多边形，那么这里涉及的是曲线的一段“无穷小”的弧，它被当作一个“无穷小”的直线段，或者是弦，或者是其存在被假定的切线上的一个线段；或者，被考虑的是一段“无穷小”的时间区间，在这期间（只要涉及的是速度）运动“是”均匀的；更大胆的是笛卡尔，他在想确定一条不服从其一般规则的摆线的切线时，把互相滚动的曲线当作多边形，以便断言“在无穷小之中”运动可以被当作绕一个接触点的旋转（{[}85 a{]}，第 II 卷，第 307–338 页）。在这里也一样，费马把他的切线法则和极大极小法则建立在这类无穷小的考虑之上，却能够在每个特定的情形中证成它们（{[}109{]}，第 I 卷，第 133–179 页；法文译本，第 III 卷，第 121–156 页）；另参看（{[}109{]}，第 II 卷，处处，特别是第 154–162 页，及《著作补遗》（Supplément aux Oeuvres）第 72–86 页）；巴罗则以古人的方式、从单调性和凸性的简单假设出发，对他定理的一大部分给出了精确的证明。但是，把新酒装进旧皮囊的时代已经过去了。在这一切之中，正如我们今天所知道的，正在成形的乃是极限的概念；而如果说，从帕斯卡、牛顿以及其他人那里，有可能提取出一些看起来与我们现代定义非常接近的陈述，那也只须把它们放回它们的语境，便能认识到那些妨碍一种严格陈述的无法移除的障碍。从 18 世纪起，当一些被清晰性需要所驱动的数学家想给他们财富的这堆乱麻理出些秩序时，这类在先辈著作中遇到的指示对他们是很珍贵的；例如，当达朗贝尔解释说，微分中除极限概念外别无他物、并精确地定义了它{[}75 b{]}时，可以相信他是受了牛顿关于“消逝量的最初与最终比”的考虑的指引{[}233 b{]}。但是，只要涉及的是 17 世纪，就必须认识到：直到牛顿和莱布尼茨背过身去抛开过去、承认必须暂时不在严格的证明中、而是在结果的多产与融贯中为他们新方法寻求证成之前，现代分析的道路是不会打通的。
\item
  动力学。我们已经看到，阿基米德早已给出他的螺线的一种动力学定义；而在中世纪，一种关于量随时间变化的初级理论（但缺乏反证，不含无穷小的考虑）以及它们的图形表示正在发展，后者的起源也许要到巴比伦天文学中去找。但对 17 世纪的数学至关重要的一点是：从一开始，微分问题就不仅联系着切线、而且联系着速度而产生。然而伽利略（{[}122 a and b{]}）在探求落体中速度的定律时（在借助斜面得到距离定律 \(x = at^2\) 之后），并不是通过微分来进行的：他对速度作出各种假设，先是 \(v = \frac{dx}{dt} = cx\)（{[}122 b{]}，第 VIII 卷，第 203 页），后来是 \(v = ct\)（\(id.\)，第 208 页），并试图借助对速度作为时间函数的图像、以相当含混的方式重新找出距离的定律；笛卡尔（1618 年）以同样的方式就定律 \(v = ct\) 进行论证，但作为一个真正的数学家，并带着不可分量语言所能汇聚的那么多清晰性\(^{5}\)（{[}85 a{]}，第 X 卷，第 75–78 页）；在这两人那里，（一条直线上的）速度图像扮演主要的角色，而且有理由问，他们对所经距离与时间轴和速度曲线所围面积之间的比例性意识到何种程度；但这一点很难说清，尽管笛卡尔的语言似乎蕴含着对所说事实的认识（某些历史学家想把这一认识追溯到中世纪{[}335{]}），而伽利略则没有对此作任何明确的评述。巴罗在 1670 年明确陈述了它（{[}16 b{]}，第 171 页）；也许它在当时对任何人都已不再是新奇之事，巴罗也没有把它当作新奇之事来呈现；但无论对这个结果还是对任何别的结果，都不宜给出过于准确的日期。至于假设 \(v = cx\)，伽利略也考虑过，他满足于（同上出处）证明它是站不住脚的（或者，用现代的语言说，方程 \(\frac{dx}{dt} = cx\) 没有一个 \(\neq 0\) 的、在 \(t = 0\) 时为零的解），用的是一段含混的论证，费马后来（{[}109{]}，第 II 卷，第 267–276 页）不辞辛劳把它展开（它差不多等于说：\(2x\) 与 \(x, x \neq 0\) 一样是解，这与解的显然的物理唯一性相悖）。但正是这同一条定律 \(\frac{dx}{dt} = cx\)，在 1614 年被纳皮尔用来引进他的对数，他对对数给出了一种动力学的定义{[}230 a{]}，用我们的记号，这定义可以写成：若在两条直线上，两个物体按定律 \(\frac{dx}{dt} = a\)、\(\frac{dy}{dt} = -ay/r\)、\(x_0 = 0\)、\(y_0 = r\) 运动，则称 \(x\) 是 \(y\) 的“对数”（用现代的记号，我们有 \(x = r \log(r/y)\)）。我们已经看到，作为 \(\frac{dy}{dx} = \frac{c}{x}\) 的解的这条曲线，于 1639 年出现在笛卡尔那里，他动力地描述了它（{[}85 a{]}，第 II 卷，第 514–517 页）；诚然，他颇为轻蔑地把一切非代数的曲线都称为“力学的”，并声称要把它们从几何学中排除出去；但幸运的是，这一禁忌——很久之后莱布尼茨还认为必须对它作激烈的抗议——既没有被他的同时代人所遵守，也没有被笛卡尔本人所遵守。摆线、对数螺线出现了，并被热情地加以研究，它们的研究是几何方法与动力学方法相互交织的一股强大助力。运动的合成的原则，更确切地说速度的合成的原则，是伽利略在他晚年的杰作、1638 年的《谈话录》（Discorsi）中所陈述的抛体运动理论的基础（{[}122 b{]}，第 VIII 卷，第 268–313 页），这一理论于是隐含地包含了抛物线切线的一种新确定；如果伽利略未曾明确陈述它，托里拆利则相反（{[}310{]}，第 III 卷，第 103–159 页）坚持这一点，甚至把这同一个原则作为对能够作动力学定义的曲线确定切线的一般方法的基础。诚然，他在这件事上被罗贝瓦尔抢了先好几年（{[}263{]}，第 3–67 页），后者说他是通过摆线的研究被引向这个方法的；正是摆线切线这同一个问题，此外还给费马提供了显示其微分方法之力量的机会（{[}109{]}，第 I 卷，第 162–165 页），而笛卡尔由于不能对它应用他的代数方法，就为这个场合发明了瞬时转动中心（{[}85 a{]}，第 II 卷，第 307–338 页）。
\end{enumerate}

但是，随着无穷小演算的发展，动力学不再是一门独立的科学。人们越来越注意到，尽管有笛卡尔，代数曲线和代数函数从无穷小演算所采取的“局部”观点看，并没有任何使它们与远为一般的函数相区别的东西；以动力学方式定义的函数和曲线是与其他函数和曲线一样的函数和曲线，同样可以为同样的方法所处理；而变量“时间”不过是一个参数，它的时间性方面纯属语言问题。于是在惠更斯那里，即使涉及的是力学，占统治地位的仍是几何（{[}169 b{]}，第 XVIII 卷）；而莱布尼茨在他的计算中不给予时间任何特殊的位置。相反，巴罗设想把若干个量作为被设想成“时间”的一个普遍独变量的函数的同时变化，做成一种带几何倾向的无穷小演算的基础。这个想法——他一定是在设法重新找出那个他只风闻其存在的运动合成方法时想到的——在他的《几何讲义》（Lectiones Geometricae）{[}16 b{]}的前三讲中，以十分清楚而十分一般的词句得到了详细的陈述；例如，他在那里细心地证明：若一个动点在两条互相垂直的轴 \(AY\)、\(AZ\) 上的投影是这样两个动点，其中一个以常速度 \(a\) 运动，另一个以随时间增长的速度 \(v\) 运动，则轨迹具有一条斜率等于 \(v/a\) 的切线，并且在 \(Z\) 增大的方向上是凸的。在《几何讲义》的后续部分中，他把这些思想推进得很远，虽然他带着某种做作、几乎自始至终以尽可能几何而尽可能少代数的形式表述它们，但人们有理由与雅各布·伯努利一起（{[}19 a{]}，第 I 卷，第 431 页及第 453 页）在那里看到牛顿和莱布尼茨无穷小演算的一大部分的等价物。正是这同样的思想构成了牛顿的出发点（{[}233 a{]}，第 I 卷，第 201–244 页及{[}233 b{]}）：他的“流量”是作为“时间”之函数的各种量，而“时间”只是一个普遍的参数，“流数”则是关于“时间”的导数；牛顿允许自己视需要更换参数，这种可能性同样存在于巴罗的方法中，尽管他运用起来不那么灵活自如。\(^{6}\) 流数的语言，为牛顿所采用、并借他的权威强加给下一个世纪的英国数学家们，就这样代表了我们所关切的时期中动力学方法的最后一座前哨，而这种方法的真正角色其实早已完结。

\begin{enumerate}
\def\labelenumi{\Alph{enumi})}
\setcounter{enumi}{2}
\item
  代数几何。这是一个寄生的主题，与我们的话题无关，它之所以在这里被引入，是因为笛卡尔以其系统化的做法，声称要使代数曲线成为几何学的唯一对象（{[}85 a{]}，第 VI 卷，第 390 页）；因此，他为确定切线所给出的，是一种代数几何的方法，而不像费马那样，是一种出自微分学的方法。古人遗下的关于直线与圆锥曲线相交的结果，笛卡尔本人关于两条圆锥曲线相交以及可化归于此的各种问题的思考，都将十分自然地把他们引向这个想法：取两个交点的重合作为相切的判据；我们今天知道，在代数几何中这是正确的判据，而且具有如此之大的普遍性，以至于它与极限的概念以及“基域”的本性都无关。笛卡尔最初以一种不甚方便的方式应用它，设法使所研究的曲线与一个圆心在 Ox 上的圆的两个交点在一个给定的点重合（{[}85 a{]}，第 VI 卷，第 413–424 页）；他的学生们范斯库滕、赫德用一条直线代替圆，并以 \(-F_{x}^{\prime}/F_{y}^{\prime}\) 的形式得到曲线 \(F(x,y)=0\) 的切线的斜率，“导出多项式”\(F_{x}^{\prime}, F_{y}^{\prime}\) 则由其形式的构成法则来定义（{[}198 d{]}，第 I 卷，第 147–344 页及{[}85 b{]}，第 234–237 页）；斯吕塞也在差不多同一时间达到这一结果（{[}198 d{]}，第 232–234 页）。当然，我们在这里所划出的这些泾渭分明的区分——唯有它们才赋予笛卡尔与费马之间的论战以意义——根本不可能存在于 17 世纪数学家们的头脑中：我们提到它们，只是为了把我们这段历史中较为奇特的一页弄得清楚些，并且紧接着指出代数方法随后的完全湮没，它被微分方法暂时吸收了。
\item
  问题的分类。我们已经看到，这个主题在阿基米德的工作中似乎是缺席的；在他看来，解决一个问题，还是把它化归为一个已经处理过的问题，仿佛是无所谓的。在 17 世纪，微分问题起初以三种不同的面貌出现：速度、切线、极大与极小。就最后一类而言，开普勒{[}179 a and b{]}注意到（这一观察在奥雷姆那里已经可以见到（{[}335{]}，第 141 页），而且甚至不曾逃过巴比伦天文学家的注意），函数的变化在一个极大值的邻域内特别缓慢。费马在 1630 年以前（{[}109{]}，第 II 卷，第 71 页及第 113–179 页）就联系这类问题开创了他的无穷小方法，用现代的语言说，这方法实际上归结为求泰勒展开中开头的两项（常数项与一阶项），并写
\end{enumerate}

出在极值处第二项为零；他由此出发，把他的方法扩展到切线的确定，甚至把它应用于拐点的寻求。如果考虑到上文就动力学所说的话，便可以看出，与一阶微分有关的三类问题的统一是相当早便完成了的。至于与二阶微分有关的问题，它们出现得很晚，而且特别是随着惠更斯关于曲线渐屈线的工作（1673 年发表于他的《摆钟论》（Horologium Oscillatorium）（{[}169 b{]}，第 XVIII 卷））；这时，牛顿以其流数早已拥有解决这类问题所必需的一切分析手段；而尽管惠更斯在其上倾注了全部几何才智（日后的微分几何在其发端时将从中受益），在我们所关切的时期内，这些问题几乎没有被用于别的什么，只是让新分析意识到自己手段的力量。

至于积分，它在希腊人那里是以面积、体积、静力矩的计算，以圆周长和球冠面积的计算的面貌出现的；17 世纪又给它添上了曲线的求长、旋转曲面面积的计算，以及（随着惠更斯关于复摆的工作（{[}169 b{]}，第 XVIII 卷））转动惯量的计算。首先的事情是认识到所有这些问题之间的联系。就面积和体积而言，这第一步的巨大进展是由卡瓦列里在他的《不可分量几何学》{[}57 a{]}中作出的。他在那里宣布并声称证明了大致如下的原则：若两个平面面积被每一条沿给定方向的平行线截出的线段其长度成固定的比，则这两个面积也成同一的比；对于被平行于一个固定平面的诸平面所截、其截面积成固定比的立体，陈述了一个类似的原则。这些原则很可能是被诸如欧几里得（或者不如说欧多克索斯）关于等高棱锥体积之比的定理提示给卡瓦列里的，而且在把它们陈述成一般的形式之前，他先用取自阿基米德的大量例子验证了它们的有效性。他借助一种语言来“证成”它们，关于这种语言的合法性，我们在他 1621 年写给伽利略的信中看到他在探问，而 1622 年他就已经毫不迟疑地使用它了（{[}122 b{]}，第 XIII 卷，第 81 页及第 86 页），其要点如下。设有例如两个面积，一个为 \(0 \leq x \leq a\)、\(0 \leq y \leq f(x)\)，另一个为 \(0 \leq x \leq a\)、\(0 \leq y \leq g(x)\)；纵标之和 \(\sum_{k=0}^{n-1} f(ka/n)\)、\(\sum_{k=0}^{n-1} g(ka/n)\) 彼此之比，当 \(n\) 充分大时任意接近于两面积之比，而且对于单调的 \(f\) 和 \(g\)，甚至不难用穷竭法证明这一点；卡瓦列里过渡到极限，取 \(n = \infty\)，并谈论第一条曲线“一切纵标之和”，它与第二条曲线相应的和之比严格地等于两面积之比；对体积也是如此；这种语言随后被普遍采用，甚至被像费马这样对它所掩盖的精确事实有最清楚认识的作者们采用。诚然，此后许多数学家，如罗贝瓦尔（{[}263{]}，第 3–67 页）和帕斯卡（{[}244{]}，第 VIII 卷，第 334–384 页），宁愿在被取“和”的曲线的纵标中看到的，不是像卡瓦列里那样的直线段，而是具有同一无穷小高度的一些矩形，从严格性的观点看这算不上多大的进步（无论罗贝瓦尔怎么说），但这也许可以防止想象太容易出轨。无论如何，既然涉及的只是比，“曲线 \(y = f(x)\) 的一切纵标之和”这个说法，或简言之曲线的“一切纵标”，最终——正如例如在帕斯卡的著作中清楚地显现的那样——是莱布尼茨的 \(\int ydx\) 的精确等价物。

从卡瓦列里所采用的语言出发，上文陈述的那些原则便不可避免地随之而来，并立即蕴含一些我们将用现代记号陈述的结论，但约定在那里 \(\int f dx\) 只表示由 \(Ox\) 与曲线 \(y = f(x)\) 所界的面积。首先，每个平面面积，只要被每条 \(x =\) 常数的直线截出的诸线段长度之和为 \(f(x)\)，就等于 \(\int f dx\)；每个立体，只要被每个 \(x =\) 常数的平面截出面积为 \(f(x)\)，也是如此。此外，\(\int f dx\) 线性地依赖于 \(f\)；我们有 \(\int (f + g)dx = \int f dx + \int g dx\)，\(\int cf(x)dx = c\int f dx\)。特别地，一切面积和体积问题都被化归为求积，也就是说化归为形如 \(\int f dx\) 的面积计算；而且，也许更为新颖也更为重要的是，依赖于同一个求积的两个问题必须看成等价的，而且在每个情形都有办法判定是否如此。希腊数学家们从未达到（或者也许从未肯于达到）这样程度的“抽象”。于是（{[}57 a{]}，第 133 页）卡瓦列里毫无困难地“证明”了两个相似立体之比等于相似比的立方，而阿基米德只在他的立体理论末尾、就旋转二次曲面及其截体陈述这一结论（{[}5 b{]}，第 I 卷，第 258 页）。但要达到那里，必须把阿基米德式的严格抛出船外。

于是就有了按与问题相关联的求积的实在或表观的困难程度、至少是暂时地对问题进行分类的手段。这正是当时的代数可用作典范之处：因为在代数中，在几何中提出的代数问题中，希腊人所关心的只是解，而 16 和 17 世纪的代数学家们开始把主要注意力转向按求解所用手段的本性对问题进行分类，从而为现代的代数扩张理论奏出了前奏；他们不仅已按问题所依赖的方程的次数对问题作了一个初步的分类，而且已经给自己提出了关于可能性的困难问题：用根式解出每一个方程的可能性（他们中有些人是不相信这一点的），等等；他们也关心把给定次数的一切问题化归为一个几何形式的类型。同样，在积分演算中，卡瓦列里的原则使他们能够立即认出，阿基米德解决的许多问题都化归为求积 \(\int x^{n}dx\)，其中 n=1,2,3；他还发明了一个巧妙的方法，可对随心所欲的多个 n 值完成这一求积（这方法归结为注意到由齐次性有 \(\int_{0}^{2a}x^{n}dx=c_{n}a^{n+1}\)，并写下

\[
\int_ {0} ^ {2 a} x ^ {n} d x = \int_ {- a} ^ {a} (a + x) ^ {n} d x = \int_ {0} ^ {a} [ (a + x) ^ {n} + (a - x) ^ {n} ] d x,
\]

展开便得 \(c_n\) 的一个递推关系）（{[}57 a{]}，第 159 页及{[}57 b{]}，第 269–273 页）。但费马那时已经走得远得多，他先（在 1636 年之前）证明了 \(\int_0^a x^n dx = \frac{a^{n+1}}{n+1}\) 对正整数 \(n\) 成立（{[}109{]}，第 II 卷，第 83 页），借助的是前 \(N\) 个整数的幂和公式（一个模仿阿基米德螺线求积而得的程序），随后把同一公式推广到一切有理的 \(n \neq 1\)（{[}109{]}，第 I 卷，第 195–198 页）；对最后这个结果（1644 年传达给卡瓦列里），他直到很久以后、在读过帕斯卡关于积分的著作\(^7\)之后，才给出一个证明（{[}109{]}，第 I 卷，第 255–288 页；法文译本第 III 卷，第 216–240 页）。

这些结果，再加上起着变量代换和分部积分作用的几何考虑，已经使大量可化归为初等求积的问题得到解决。在此之外，首先遇到的是圆的求积和双曲线的求积：由于那个时代的人特别关心的是“不定积分”，这些问题用现代的术语说，其解分别由反三角函数和对数提供；这些解是以几何方式给出的，而我们也已经看到，这后一个函数是怎样一步一步被引进分析的。这些求积成了大量工作的对象，从事者有格雷古瓦·德·圣万桑{[}135{]}、惠更斯（{[}169 b{]}，第 XI 卷，第 271–337 页及第 XII 卷，第 91–180 页）、沃利斯（{[}326 a{]}，第 I 卷，第 355–478 页）、格雷戈里{[}136 a{]}；第一位自信已经完成了圆的求积，最后一位自信已经证明了 \(\pi\) 的超越性；在他们那里，为三角函数和对数函数发展出了不定的逼近程序，一些偏于理论，另一些面向数值计算，而这些程序很快就要以牛顿（{[}233 a{]}，第 I 卷，第 3–26 页及第 29–199 页）、墨卡托（{[}220{]}，转载于{[}218{]}，第 I 卷，第 167–196 页）、J. 格雷戈里{[}136 d{]}、然后是莱布尼茨{[}198 d{]}的级数展开的一般方法而告终。无论如何，人们渐渐确信这些求积的“不可能性”，也就是说，确信它们所定义的函数的非代数特征；与此同时，形成了一种习惯：一个问题，当它已被化归为“不可能的”求积之一时，就算就其本性所允许的程度而言解决了。摆线的问题便是如此，它借助三角函数得到解决，抛物线的求长也是如此，它被化归为双曲线的求积。

求长问题——我们刚才引了其中最著名的两个——具有一种特殊的重要性，因为它们构成微分（它们以之为前提）与积分（它们所属范围）之间的一个自然的几何过渡；旋转曲面面积的问题可以与它们相联系。古人只处理过圆和球的情形。在 17 世纪，这些问题很晚才出现；看来是椭圆求长（椭圆被视为圆之后最简单的曲线）的、该时代无法逾越的困难使人望而却步了。动力学方法给这些问题提供了某些入口，使罗贝瓦尔（{[}263{]}，第 363–399 页）和托里拆利（{[}310{]}，第 III 卷，第 103–159 页）在 1640 年到 1645 年之间得到了关于螺线弧的一些结果；但只是到了 1660 年前的那些年，它们才成为当务之急；摆线于 1658 年被雷恩求长（{[}326 a{]}，第 I 卷，第 533–541 页）；不久之后，曲线 \(y^3 = ax^2\) 被好几位作者处理（{[}326 a{]}，第 I 卷，第 551–553 页；{[}85 b{]}，第 I 卷，第 517–520 页；{[}109{]}，第 I 卷，第 211–255 页；法文译本，第 III 卷，第 181–215 页），还有好几位作者（{[}109{]}，第 I 卷，第 199 页；{[}169 b{]}，第 II 卷，第 334 页）把抛物线的求长化归为双曲线的求积（也就是说化归为一个代数–对数函数）。最后这个例子最为重要，因为它是下述一般原理的一个特例：曲线 \(y = f(x)\) 的求长不是别的，正是 \(y = \sqrt{1 + (f'(x))^2}\) 的求积；而惠拉特正是从这条原理本身把它推出来的。跟踪费马晚年就曲线 \(y^3 = ax^2\) 的工作（{[}109{]}，第 I 卷，第 211–255 页；法文译本第 III 卷，第 181–215 页）中的摸索也同样有趣：对曲线 \(y = f(x)\) 及其弧 \(s = g(x)\)，他联结曲线 \(y = g(x)\)，并从第一条曲线的切线出发确定这一条曲线的切线（用现代的语言说，他证明它们的斜率 \(f'(x), g'(x)\) 由关系 \((g'(x))^2 = 1 + (f'(x))^2\) 联系）；人们会觉得自己离巴罗非常之近，只须把这个结果与惠拉特的结果结合起来（这正是格雷戈里 1668 年差不多做到的（{[}136 d{]}，第 488–491 页）），便可得到切线与求积之间的关系；但费马陈述的只是：若两条曲线各自相对于一个直角坐标系，在具有相同横标的点处的切线总有相同的斜率，则这两条曲线相等，或者换句话说，\(f'(x)\) 的知识确定 \(f(x)\)（至多差一个常数）；而他对这个断言的证成，只是借助一个没有任何说服价值的含混论证。

不到十年之后，巴罗的《几何讲义》{[}16 b{]}问世了。从一开头（Lect. I）他就以原理的形式陈述：在直线运动中，诸空间与由时间轴和速度曲线所界的面积 \(\int_{0}^{t} v dt\) 成比例。人们会以为，他将由此、并从他先前已引述过的关于切线确定的动力学方法，推出被视为切线斜率的微分与被视为面积的积分之间的联系；但事情并非如此，他后来以纯几何的方式证明（{[}16 b{]}，Lect. X, §11, 第 243 页）：若两条曲线 \(y = f(x), Y = F(x)\) 满足纵标 \(Y\) 与面积 \(\int_{a}^{x} y dx\) 成比例，也就是说若 \(c.F(x) = \int_{a}^{x} f(x) dx\)，则 \(Y = F(x)\) 的切线与 \(Ox\) 相交于横标为 \(x - T\) 的点，该横标由 \(y/Y = c/T\) 确定；这个证明此外是完全精确的，从 \(f(x)\) 单调这一显式假设出发，而且书中说到，\(f(x)\) 的变化方向确定 \(Y = F(x)\) 的凹向。但将可注意到，这条定理多少淹没在一大堆别的定理之中，其中许多是非常有趣的；未曾预先受到提醒的读者会倾向于在那里只看到用求积解决问题 \(Y/T = f(x)/c\) 的一个手段，也就是说，看到从切线信息出发确定一条曲线的某个问题（或者如我们所说，某种类型的微分方程）；而既然巴罗所给出的它的应用尤其涉及同类的问题（也就是说涉及可用“分离变量法”积分的微分方程），就更是如此。巴罗强加于自己的几何语言在这里是一个原因，它使微分与积分之间的联系——只要涉及的是动力学，这联系原本那么清楚——多少被遮蔽了。

另一方面，各种方法已经成形，用来把一些积分问题化归为另一些，并“解出”它们，或把它们化归为已经分类的“不可能”问题。分部积分以其最简单的几何形式，就在于把由 \(Ox, Oy\) 和单调曲线 \(y = f(x)\) 的一段弧（这段弧把 \((a, 0)\)（\(Ox\) 上的点）与 \((0, b)\)（\(Oy\) 上的点）连接起来）所界的面积写成 \(\int_{0}^{a} ydx = \int_{0}^{b} xdy\)；而且它经常以隐含的方式被使用。随着帕斯卡（{[}244{]}，第 IX 卷，第 17–18 页）出现了如下的推广，它隐藏得深得多：设 \(f(x)\) 如上，\(g(x)\) 是一个 \(\geq 0\) 的函数，且 \(G(x) = \int_{0}^{x} g(x)dx\)；则有 \(\int_{0}^{a} yg(x)dx = \int_{0}^{b} G(x)dy\)，他巧妙地通过用两种方式计算立体 \(0 \leq x \leq a, 0 \leq y \leq f(x), 0 \leq z \leq g(x)\) 的体积证明了它；特例 \(g(x) = x^n\)，\(G(x) = \frac{x^{n+1}}{n+1}\) 起着重要的作用，在帕斯卡（同上出处，第 19–21 页）那里和在费马（{[}109{]}，第 I 卷，第 271 页）那里一样；这后一位（其著作带着一个意味深长的标题：《曲线方程的变换与订正，及其对曲线空间相互之间以及与直线空间之比较的种种应用\ldots{}》）没有证明它，无疑是因为他觉得重复帕斯卡刚刚发表的东西并无用处。这些“变换”定理——我们在其中会看到分部积分与变量代换的一种结合——在某种程度上代替了它，后者是很晚才被引入的；事实上，与当时的思想方式——仍然过于几何、不够分析——相抵触的，是允许使用图形所强加的变量以外的变量，也就是说两坐标之一（有时是极坐标）或曲线的弧。正是这样，我们在帕斯卡（{[}244{]}，第 IX 卷，第 60–76 页）那里找到一些结果，用现代的记号写出、令 \(x = \cos t\)、\(y = \sin t\)、并对特殊的函数 \(f(x)\)，就是

\[
\int_ {0} ^ {1} f (x) d x = \int_ {0} ^ {\frac {\pi}{2}} f (x) y d t,
\]

而与 J. 格雷戈里（{[}136 d{]}，第 489 页）那里相当，对一条曲线 \(y = f(x)\) 及其弧 \(s, \int yds = \int zdx\)，这里 \(z = y\sqrt{1 + y'^2}\)。只是到了 1669 年，我们才看到巴罗掌握了变量代换的一般定理（{[}16 b{]}，第 298–299 页）；它的陈述一如既往是几何的，可归结如下：设 \(x\) 与 \(y\) 由一个单调关系相联系，\(p\) 是该关系的图像在点 \((x, y)\) 处的斜率；那么，若函数 \(f(x), g(y)\) 使得 \(f(x)/g(y) = p\) 对每一对对应的值 \((x, y)\) 都成立，则在相应的限之间取的面积 \(\int f(x)dx, \int g(y)dy\) 相等；而且反过来，若这些面积总是相等（\(f\) 和 \(g\) 隐含地被设为具有常号），则我们有 \(p = f(x)/g(y)\)；这个逆命题自然被用来把定理应用于微分方程的求解（用“分离变量法”）。但这条定理只是被巴罗插在一个附录中（Lect. XII, app. III, theor. IV），他在那里一面指出他先前的结果中有几条只是它的特例，一面对自己发现它太迟、来不及更多地使用它表示歉意。

于是，1670 年前后，局势如下。人们已经知道如何用统一的方法处理涉及一阶导数的问题，而惠更斯已经着手处理涉及二阶导数的几何问题。人们知道如何把一切积分问题化归为求积；已经有多种带有几何面貌的办法，可以在分类不善的情形下把一些求积化归为另一些求积，而且从这个观点出发，处理三角函数和对数函数已经习以为常；人们已经意识到微分与积分之间的联系；“切线反方法”——当时给予化归为一阶微分方程的问题的名称——也开始被着手处理。墨卡托对级数 \(\log(1+x) = -\sum_{1}^{\infty}(-x)^{n}/n\) 的轰动性发现，刚刚为把级数、且主要是幂级数应用于所谓“不可能”的问题开辟了全新的前景。与之相对照，数学家的队伍明显变稀了：巴罗离开了教授的席位去就任布道者的席位；除惠更斯之外（他已把几乎全部数学工作留在身后，早已获得他正着手最终写定的《摆钟论》的一切主要结果），只有剑桥的牛顿和阿伯丁的 J. 格雷戈里还在活跃，而莱布尼茨不久就将以新入教者的热忱加入他们。这三个人，牛顿从 1665 年起、J. 格雷戈里从 1668 年墨卡托的发表起、莱布尼茨从约 1673 年起，都主要投身于当日的时髦课题，即幂级数的研究。但是，从问题分类的观点看，新方法的主要效果看来是抹掉了它们之间的一切区分；而实际上，牛顿——与其说他是代数学家，不如说他是分析学家——在 1676 年毫不犹豫地向莱布尼茨宣布（{[}198 d{]}，第 224 页），他知道如何解一切微分方程，\(^{8}\) 对此莱布尼茨回答（{[}198 d{]}，第 248–249 页）说，恰恰相反，这是一个每当可能时就要“预设求积”而以有限数目的项得出解的问题，也是一个要知道是否每个求积都能像迄今研究的多数情形中已被注意到的那样，化归为圆和双曲线的求积的问题；他借此重新提出，格雷戈里相信（有理由，正如我们今天所知）椭圆和双曲线的求长不可化归为圆和双曲线的求积；莱布尼茨并问道，级数的方法，如牛顿所使用的那种，能在多大程度上对这些问题给出回答。牛顿在他那方面（{[}198 d{]}，第 209–211 页）声称他握有一些他不予指明的判据，可用来——表面上是通过考察级数——判定某些（以有限项表示的）求积的“可能性”，并给出了积分 \(\int x^{\alpha}(1 + x)^{\beta} dx\) 的一个（非常有趣的）级数的例子。

十年之内实现的巨大进步由此可以看出；分类问题在这些通信中是以十分现代的术语陈述的；莱布尼茨提出的问题之一已在 19 世纪由阿贝尔积分的理论解决，另一个，即把一个给定的微分方程化归为求积的可能性问题，尽管近来有重要的工作，仍然悬而未决。如果事情就是这样，那是因为牛顿和莱布尼茨已经各自把无穷小演算的基本运算化归为一个算法；只须用两人中这个或那个所使用的记号写下一个关于求积或微分方程的问题，它的代数结构便会立即显现出来，摆脱了它的几何外壳；“变换”的方法同样可以用简单的分析术语写出；分类问题得到了精确的陈述。从数学上讲，17 世纪已经达到了它的终点。

\begin{enumerate}
\def\labelenumi{\Alph{enumi})}
\setcounter{enumi}{4}
\tightlist
\item
  内插法与差分演算。这个主题（我们不把二项式系数的研究与它分开）出现得很早，并贯穿整个世纪，其理由既有理论上的也有实践上的。当时的一大任务事实上是三角函数表、对数表、航海表的计算，地理学、航海术、理论和实用天文学、物理学、天体力学的迅速进展使它们成为必需；许多最杰出的数学家，从开普勒到惠更斯和牛顿，都参与了这项工作，或者是直接参与，或者是通过关于最有效逼近程序的理论研究参与。
\end{enumerate}

在数表的使用乃至编制中，最初的问题之一就是内插；而随着计算的精度增长，17 世纪注意到，古老的线性内插程序，在最初差分（表中相继出现的数值之间的差）不再实际上是常数时，便失去了效力；例如，我们看到布里格斯（{[}36{]}，第 13 章）在对数的计算中使用高阶的差分、甚至是相当高阶的差分。后来，我们看到牛顿（{[}233 a{]}，第 I 卷，第 271–282 页及{[}233 b{]}，第 III 卷，引理 5；亦见{[}114{]}）和 J. 格雷戈里（{[}136 d{]}，第 119–120 页）各自独立地进行着关于内插和幂级数的平行研究；两人以此外并不相同的方法，一方面得到多项式内插公式，即所谓“牛顿的”公式，另一方面得到二项级数（{[}136 d{]}，第 131 页；{[}198 d{]}，第 180 页）和经典分析中主要的级数展开（{[}136 d{]}；{[}233 a{]}，第 I 卷，第 3–26 页及第 271–282 页，及{[}198 d{]}，第 179–192 页及第 203–225 页）；这两条研究路线彼此激荡，而且在牛顿心中与无穷小演算诸原理的发现紧密相连，这几乎是不容置疑的。在格雷戈里那里如同在牛顿那里一样，一种对数值实践的巨大关切显露出来：数表的编制与使用、级数与积分的数值计算；特别地，虽然他们的著作中找不到一个像上引布龙克尔勋爵那样的打磨过的收敛性证明，两人都不断地从他们的级数之于计算的适用性这一实践观点谈到级数的收敛。我们仍看到牛顿这样行事：为回答柯林斯出于实用目的提出的一个问题（{[}262{]}，第 II 卷，第 309–310 页），他把所谓欧拉–麦克劳林求和方法的一个特例应用于 \(\sum_{p=1}^{N} \frac{1}{n+p}\) 当 \(N\) 值很大时的近似计算。

从差分出发计算一个函数的值，也早早地遇到，被用作积分的一个实用程序，甚至可以说，被用作微分方程积分的实用程序。例如赖特在 1599 年，为编制航海表而必须解一个我们应记作 \(\frac{dx}{dt} = \sec t = \frac{1}{\cos t}\) 的问题时，按一弧秒的相继间隔把 \(\sec t\) 的值逐一相加（{[}338{]}；参见{[}230 b{]}，第 97 页），自然实际上得到了 \(\log \tan (\frac{\pi}{4} + \frac{t}{2})\) 的数值表；而这一从最初几张 \(\log \tan t\) 数表的计算中观察到的巧合，直到 1668 年格雷戈里做出 \(\sec t\) 的积分之前，一直没有得到解释（{[}136 c{]}及{[}136 d{]}，第 7 页及第 463 页）。

但这些问题也有一个纯理论的、甚至算术的方面。我们约定以 \(\Delta^r x_n\) 记一个序列 \((x_n)_{n\in \mathbf{N}}\) 的相继差分，它借助 \(\Delta x_{n} = x_{n + 1} - x_{n}\)、\(\Delta^r x_n = \Delta (\Delta^{r - 1}x_n)\) 递归地定义，并以 \(S^r\) 记 \(\Delta\) 及其迭代的逆运算，也就是说，若 \(y_{n} = Sx_{n}\)，其中 \(y_0 = 0\)、\(\Delta y_{n} = x_{n}\)，且 \(S^r x_n = S(S^{r - 1}x_n)\)，便作此约定；我们有

\[
S ^ {r} x _ {n} = \sum_ {p = 1} ^ {n - r} \binom{n - p - 1}{r - 1} x _ {p};
\]

特别地，若对一切 n 有 \(x_{n}=1\)，则 \(Sx_{n}=n\)，序列 \(S^{2}x_{n}\) 和 \(S^{3}x_{n}\) 就是希腊算术家已经研究过的“三角形数”和“角锥数”的序列，而一般地我们有

\[
S ^ {r} x _ {n} = \binom{n}{r} \text { for } n \geq r (\text { and } S ^ {r} x _ {n} = 0 \text { for } n <   r);
\]

这些序列从这个观点看至少从 16 世纪起就被引入了；它们也自行出现在组合问题中，这些问题无论是凭其自身、还是联系着概率，都在 17 世纪的数学中起了相当大的作用，例如在费马和帕斯卡那里，然后在莱布尼茨那里。它们也出现在 \(m\) 次幂之和的表达式中，这个和是对前 \(N\) 个整数取的，这个和的计算，正如我们已经看到的，是费马以第一种方法积分 \(\int x^{m}dx\)（\(m\) 为整数）的基础（{[}109{]}，第 II 卷，第 83 页）。沃利斯 1655 年在他的《无穷算术》（Arithmetica Infinitorum）（{[}326 a{]}，第 I 卷，第 355–478 页）中正是这样进行的，他既不知道费马的（未发表的）工作，而且他说，对不可分量法，他也只是从阅读托里拆利中有所闻；诚然，急于到达终点的沃利斯不让自己被细致的研究拖住脚步：一旦对 \(m\) 的前几个整数值得到了结果，他便“用归纳”断言它对一切整数 \(m\) 为真，从这里正确地过渡到 \(m = 1/n\)（\(n\) 为整数），然后以一种比第一次更加简略的“归纳”过渡到任意的有理数 \(m\)。但他工作中有趣而新颖之处在于，他从这里逐步上升到对“欧拉的”积分 \(I(m,n) = \int_{0}^{1}(1 - x^{\frac{1}{m}})^{n}dx\)（其值当 \(m\) 和 \(n > 0\) 时为 \(\Gamma(m+1)\Gamma(n+1)/\Gamma(m+n+1)\)）以及别的一些这样的积分的研究，对整数 \(m\) 和 \(n\) 编制 \(1/I(m,n)\) 的数值表，它不是别的，正是整数 \(\binom{m+n}{n}\) 的那张表，而且，他以与今天讲述 \(\Gamma\) 函数理论所用几乎相同的方法，最终得到了 \(I\left(\frac{1}{2},\frac{1}{2}\right)=\frac{\pi}{4}=\left(\Gamma\left(\frac{3}{2}\right)\right)^{2}\) 的无穷乘积；此外，借助分部积分和非常简单的变量代换、并通过把 \(I(m,n)\) 就 \(m\) 和 \(n\) 的一切实数值加以考虑，把这方法搞正确并不困难，只是这一点他几乎不可能想到，而牛顿的分析不久就将使它成为可能。无论如何，正是沃利斯对整数 \(\binom{m+n}{n}\) 向 \(m\) 的非整数值（更确切地说，向形如 \(n=p/2\)（\(p\) 为奇整数）的值）所作的“内插”，为起步时的牛顿（{[}198 d{]}，第 204–206 页）提供了一个出发点，把他先引向特例 \((1-x^{2})^{p/2}\) 的研究，再引向二项级数，从那里引向对 \(x^{a}\)（如此记写）就一切实数 \(a\) 的引入，并引向借助二项级数对 \(x^{a}\) 的微分；这一切都没有为得到严格的证明、乃至定义付出多大的努力；此外，一项值得注意的革新是：他是从 \(x^{a}\) 的导数的知识推出 \(\int x^{a}dx\)（\(a \neq -1\)）的（{[}233 a{]}，第 I 卷，第 3–26 页及{[}198 d{]}，第

225 页）。至于其余的，而且尽管他不久就掌握了远为一般的幂级数展开方法，例如所谓牛顿多边形的方法（用于代数函数）（{[}198 d{]}，第 221 页）和待定系数法，他在随后的工作中仍常常带着某种偏爱回到二项级数及其种种推广；而且正是从那里，例如，他似乎得出了上文所引的 \(\int x^{\alpha}(1 + x)^{\beta} dx\) 的展开（{[}198 d{]}，第 209 页）。

与此同时，大陆上思想的演变很不一样，而且抽象得多。帕斯卡在二项式系数的研究上已与费马接近（他从这些系数构造出他所说的“算术三角形”），并研究了它们在概率计算和差分计算中的用途；当他着手积分时，他把同样的思想引到了那里。像他的先辈们一样，当他使用不可分量的语言时，他把积分 \(F(x) = \int_{0}^{x} f(x) dx\) 设想为“曲线的一切纵标之和”

\[
S \left(f \left(\frac {n}{N}\right)\right) = \sum_ {0 \leq p <   N x} f \left(\frac {p}{N}\right),
\]

对“单位”\(N = \sum_{0 \leq p < N} 1\) 当 \(N = \infty\) 时的比值（{[}244{]}，第 VIII 卷，第 352–355 页）（或者，当他抛弃这种语言而采用正确的穷竭语言时，作为该比值当 N 无限增长时的极限）。但是，由于心目中惦念着矩量问题，他注意到，当涉及的是按等距间隔分布的离散质量 \(y_{i}\) 时，总质量的计算化归为上文定义的运算 \(Sy_{n}\)，而矩量的计算化归为运算 \(S^{2}y_{n}\)；而且通过类比，他把运算 \(\int\) 迭代，构成他所说的“纵标的三角和”，用我们的语言说，就是诸和 \(N^{-2}S^{2}\left(f\left(\frac{n}{N}\right)\right)\) 当 N 趋于无穷时的极限，也就是说积分 \(F_{2}(x) = \int_{0}^{x} F(x)dx\)；新的一次迭代给他“角锥和”\(F_{3}(x) = \int_{0}^{x} F_{2}(x)dx\)，即 \(N^{-3}S^{3}\left(f\left(\frac{n}{N}\right)\right)\) 的极限；此外，上下文足以表明，如果他到此为止，那既不是由于他的思想缺乏普遍性，也不是由于他的语言缺乏普遍性，而只是因为他只打算使用这些和，其系统的运用是他一大部分结果的基础，而且他立即证明了我们会写作

\[
F _ {2} (x) = \int_ {0} ^ {x} (x - u) f (u) d u, F _ {3} = \frac {1}{2} \int_ {0} ^ {x} (x - u) ^ {2} f (u) d u
\]

（{[}244{]}，第 VIII 卷，第 361–367 页）的那些性质；这一切都没有写下一个公式，而是用一种如此透明而如此精确的语言，以致人们可以立即把它转录成公式，就像刚才所做的那样。在帕斯卡那里如同在他的先辈们那里一样，但以清楚得多也系统得多的方式，独变量的选择（它总是坐标之一，或曲线的弧）隐含在确定积分区间诸等距分点（尽管“无穷接近”）的约定之中；这些点分别或者在 Ox 上，或者在 Oy 上，或者在曲线的弧上，而帕斯卡留意不让任何含混在此留存（{[}244{]}，第 VIII 卷，第 368–369 页）。当他必须更换变量时，他借助一条原则来做，这条原则可归结为说：面积 \(\int f(x)dx\) 可以写成 \(S(f(x_i)\Delta x_i)\)，对积分区间分成“无穷小”区间 \(\Delta x_i\)（相等或不等）的每一个分划都成立（{[}244{]}，第 IX 卷，第 61–68 页）。

由此可见，我们已经离莱布尼茨非常近了；而且可以说是一件幸事：后者当他想要被引进现代数学之门时，遇到的是惠更斯，后者立即把帕斯卡的著作放到了他手里（{[}198 d{]}，第 407–408 页）；他对组合分析的思考使他对这方面作了特别好的准备，而且我们知道，他对这些著作下过深入的功夫，这反映在他的工作之中。1675 年，我们看到他把上文帕斯卡的定理转录成 \(\mathrm{omn}(x\omega) = x . \mathrm{omn}\omega - \mathrm{omn}(\mathrm{omn}\omega)\) 的形式，其中 \(\mathrm{omn}\omega\) 是 \(\omega\) 的从 0 取到 \(x\) 的积分的一个缩写，几天之后，莱布尼茨用 \(\int \omega\)（“summa omnium \(\omega\)”，即“总和”的首字母）代替了它，同时引入 \(d\) 表示无穷小的“差”，或者如他不久将说的，微分（{[}198 d{]}，第 147–167 页）。他把这些“差”设想为彼此之间可以比较、却不可以与有限量比较的量，此外他最经常地、或明或暗地取微分 \(dx\)（独立变量 \(x\) 的微分）为单位，\(dx = 1\)（这相当于把微分 \(dy\) 与导数 \(\frac{dy}{dx}\) 等同起来），起先还把它从积分的记号中省去，于是积分表现为 \(\int y\) 而不是 \(\int ydx\)；但他不久就引入了它，并且一旦认识到它相对于独变量选择的不变特征，就系统地坚持它，这使他免于把这个选择一直惦记在心；\(^{9}\) 而当他回到他直到那时一直忽视的巴罗的研究时，他带着不小的得意注意到，巴罗如此看重的一般变量代换定理，从他自己 的记号出发立即就能得出（{[}198 d{]}，第 412 页）。在这一切中，他实际上始终紧贴着差分演算，他的微分演算可以由差分演算通过过渡到极限得出，当然这个过渡他很难严格地证成；而且在随后的工作中，他自愿地强调他的原则不加改变地适用于两者。他明确地引用帕斯卡，例如，在与约翰·伯努利的通信（{[}198 a{]}，第 III 卷，第 156 页）中，提到自己最初的研究时，他给出一个计算差分的公式，它是牛顿公式的一个特例，并由它“过渡到极限”得出公式 \(y = \sum_{1}^{\infty} (-1)^{n} \frac{d^{n}y}{dx^{n}} \cdot \frac{x^{n}}{n!}\)（其中 \(y\) 是一个在 \(x = 0\) 时为零的函数，而诸 \(\frac{d^{n}y}{dx^{n}}\) 是它在变量的值 \(x\) 处的导数），这个公式等价于伯努利传达给他的一个类似的公式（{[}198 a{]}，第 III 卷，第 150 页及{[}20 a{]}，第 I 卷，第 125–128 页），后者以逐次分部积分证明了它。这个公式，可以看出，非常接近泰勒级数；而莱布尼茨由差分演算出发过渡到极限的那个论证，正是泰勒 1715 年为得到“他的”级数{[}305{]}所重新找到的，\(^{10}\) 此外他对它也没有多加使用。

\begin{enumerate}
\def\labelenumi{\Alph{enumi})}
\setcounter{enumi}{5}
\tightlist
\item
  人们想必已经在上述演进中隐含地看到无穷小分析的逐步代数化，也就是说，看到它被归结为一种运算微积，配有具有代数特征的统一记号体系。正如莱布尼茨曾以完全的明晰性多次指出的那样（{[}198 a{]}，第 V 卷，第 230-233 页），这里的问题是，为新分析去做韦达曾为方程论做过的事、笛卡尔曾为几何学做过的事。要理解这种必要性，只需读几页巴罗的著作；任何时候都不可能眼前不摆着一幅图形，有时相当复杂，并且事先经过一丝不苟的精心描绘；构成这部著作实质部分的那 100 页（Lect. V-XII）竟需要不少于 180 幅图形。
\end{enumerate}

诚然，在几何表象的多样性之中出现某种统一性之前，是谈不上代数化的。尽管如此，格雷古瓦·德·圣万桑（{[}135{]}）已经（以“ductus plani in planum”之名）引进了一种合成律，它归结为对积分 \(\int_{a}^{b} f(x) g(x) dx\) 的系统使用，这些积分被看作立体 \(a \leq x \leq b, 0 \leq y \leq f(x), 0 \leq z \leq g(x)\) 的体积；但他远没有从中引出帕斯卡后来正如我们所见的那样从同一立体的研究中推出的那些结论。沃利斯在 1655 年、帕斯卡在 1658 年，各自为自己的用途锻造了具有代数特征的语言，他们不写任何公式，却拟定出一些命题，只要理解了其机制，便可立即把这些命题改写为积分演算的公式。帕斯卡的语言特别清晰而精确；如果说人们无法理解他为什么拒绝让自己使用代数记号——不仅拒绝笛卡尔的，甚至连韦达的也不用——那么人们只能叹服他所完成的这一杰作，而他对语言的驾驭乃是他做到这一点的唯一手段。

但是只要过去几年，一切就都变了。牛顿首先想出一个主意：把当时无穷小分析中一切具有几何特征的运算，都用一种唯一的分析运算即微分、加上逆问题的求解来代替；而人们很清楚，幂级数方法使他得以极其顺手地完成这种运算。正如我们已经看到的，他的语言借自于一个普适“时间”参数的虚构：他把随这个参数而变化的量称为“流量”，把它们的导数称为“流数”。他似乎没有赋予记号以特别的重要性，他忠实的崇拜者后来甚至夸口说，缺少系统的记号反倒是一种优点；不过他很早就养成了为个人所用以一点来记流数的习惯，即把 \(\frac{dx}{dt}\) 写作 \(\dot{x}\)，把 \(\frac{d^{2}x}{dt^{2}}\) 写作 \(\ddot{x}\)，等等。至于积分，牛顿看来恰似巴罗一样，从来不把它看作别的，只把它看作一个问题（已知流数求流量，即解 \(\dot{x}=f(t)\)），而不看作一种运算；因此他没有给积分的名称，似乎也没有标准的记号（只是有时用一个方框 \(\boxed{f(t)}\) 或 \(\Box f(t)\) 来表示 \(\int f(t)dt\)）。这会不会是因为，他不喜欢给一个并非唯一确定、而只确定到相差一个加法常数的对象以名称和符号？）由于缺乏文献，这个问题只能被提出来而已。

牛顿有多经验、具体、即便在最大胆处也多审慎，莱布尼茨就有多系统化、多概括化、多敢于冒险乃至有时多鲁莽。他从青年时代起就怀着一种“特征术”、即普适符号语言的想法，这种语言之于全部人类思想，将如代数记号之于代数，其中每个名称或符号都是被符号化对象的一切性质的钥匙，而且若不同时正确地推理，就不可能正确地使用它（见第 6 页）。把这样一个计划当作空想是容易的；但并非出于偶然，这个计划的作者恰恰就是那个人：他不久便认识并分离出无穷小演算的基本概念，并以多少具有最终形态的记号把它装备起来。我们已经在上文目睹了这些记号的诞生，并注意到莱布尼茨——他似乎意识到自己的使命——以怎样的细心逐步修改它们，直到确保他所寻求的简洁性尤其是不变性（{[}198 a{]}，第 V 卷，第 220-226 和 226-233 页）。这里需要注意的重要之点是：他一引进这些记号（那时他对牛顿的思想还一无所知），就对 \(\int\) 和 \(d\)、即积分和微分，有了清楚的构想，把它们视为互逆的算子。诚然，这样做他就无法避免不定积分所固有的歧义，这是他的体系的弱点，他像他的后继者一样巧妙地绕过了它。但从新符号初次出现起，引人注目的是，莱布尼茨立即忙于制定使用规则，自问是否 \(d(xy) = dxdy\)（{[}198 d{]}，第 165-166 页），并自行给出否定的回答，从而逐步得到正确的法则（{[}198 a{]}，第 V 卷，第 220-226 页），后来他又用他关于 \(d^n(xy)\) 的著名法则把这一法则加以推广（{[}198 a{]}，第 III 卷，第 175 页）。当然，在莱布尼茨如此摸索之时，牛顿早已知道 \(z = xy\) 蕴含 \(\dot{z} = \dot{x}y + x\dot{y}\) 有十年之久了；但他从不费心说出这一点，只把它看作他的流量间关系 \(F(x,y,z) = 0\) 求微分法则的一个不配命名的特例。相反，莱布尼茨的主要关切，不是使他的方法对解决这类具体问题有用，甚至也不是从严格而无可指摘的原理推出它们，而首先是建立起一种以少数简单规则的形式操纵为根基的算法。正是在这种精神下，他以括号的用法改进了代数记号，逐步采用 \(\log x\) 或 lx 来记对数，\(^{11}\) 并强调“指数演算”，也就是对指数 \(a^{x}, x^{x}, x^{y}\)（指数为变量）的系统考虑。尤其是，当牛顿只是严格按照每个具体情形的需要引入高阶流数时，莱布尼茨却很早就转向通过 d 与 \(\int\) 的迭代来创立一种“运算微积”；逐渐意识到他的演算中数的乘法与算子的复合之间的类比，他以一种幸运的大胆采用指数记号来书写 d 的迭代，把第 n 次迭代写作 \(d^{n}\)（{[}198 d{]}，第 595 和 601 页，\(^{12}\) 以及 {[}198 a{]}，第 V 卷，第 221 和 378 页），甚至用 \(d^{-1}, d^{-n}\) 表示 \(\int\) 及其迭代（{[}198 a{]}，第 III 卷，第 167 页）；他甚至寻求给 \(d^{\alpha}\) 以意义，其中 \(\alpha\) 为任意实数（{[}198 a{]}，第 II 卷，第 301-302 页，以及第 III 卷，第 228 页）。

这并不是说莱布尼茨不也关心他的演算的应用，他深知（正如惠更斯常对他反复说的那样（{[}198 d{]}，第 599 页））这些应用是它的试金石；但他缺乏深入其中的耐心，并在其中首先寻求提出新的一般规则的机会。就这样，在 1686 年（{[}198 a{]}，第 VII 卷，第 326-329 页）他处理了曲线的曲率和密切圆，并在 1692 年（{[}198 a{]}，第 VII 卷，第 331-337 页）以关于平面曲线接触的一般原理告终；\(^{13}\)在 1692 年（{[}198 a{]}，第 V 卷，第 266-269 页）和 1694 年（{[}198 a{]}，第 V 卷，第 301-306 页），他奠定了包络理论的基础；与约翰·伯努利同时，他在 1702 年和 1703 年以分解为简单元素的方式实现了有理分式的积分，但起初纯粹是形式地做的，并未充分认识到分母中出现复线性因子时伴随的情况（{[}198 a{]}，第 V 卷，第 350-366 页）。又是这样，1697 年 8 月的一天，他在马车里思考变分演算的问题时，想到了在记号 \(\int\) 之下关于一个参数求微分的法则；他热情高涨，立即把它寄给了伯努利（{[}198 a{]}，第 III 卷，第 449-454 页）。但当他做到这一步时，他的演算的基本原理早已确立，其使用已开始传播：无穷小分析的代数化已是既成事实。

\begin{enumerate}
\def\labelenumi{\Alph{enumi})}
\setcounter{enumi}{6}
\tightlist
\item
  函数的概念在 17 世纪期间以多种多样的方式被引进和精确化。全部动力学都建立在一个关于随时间变化的量、也就是时间的函数的直观的、在某种意义上实验性的观念之上，而我们已经看到，由此如何得到一个单参数的函数，正如它在巴罗那里、以及以流量之名在牛顿那里所显现的那样。“任意曲线”的观念经常出现，但很少被精确化；很可能它常被以动态的形式、或无论如何以实验的形式来设想，而且人们并不认为一条曲线要能用作推理的对象就必须具有几何的或解析的刻画；情形尤其如此（其理由我们今天更能理解）当涉及积分时，例如在卡瓦列里、帕斯卡和巴罗那里；这最后一位在推理中处理由 x = ct、\(y = f(t)\) 定义的曲线并假设 \(\frac{dy}{dt}\) 递增时，甚至明确地说，\(\frac{dy}{dt}\) “按任意规律正规地递增、甚至不正规地递增”都“无关紧要”（{[}16 b{]}，第 191 页），也就是说，正如我们会说的那样，是否能有一个解析定义。可惜这个清楚而富有成果的观念——它一经适当精确化便将在 19 世纪重新出现——未能战胜笛卡尔造成的混乱：后者一方面把一切不能精确定义为解析式的曲线从“几何学”中驱逐出去，另一方面又把这样一种定义中可容许的形成手续仅限于代数运算。诚然，在最后这一点上，他的大多数同时代人并没有跟随他；渐渐地，而且往往经由十分迂回的路径，各种超越运算——对数、指数、三角函数、求积、微分方程的求解、取极限、级数求和——获得了被引用的资格，而在每种情形都难以标出迈出这一步的确切时刻；况且第一步之后往往跟着倒退的一步。例如就对数而言，应当视为重要阶段的有：对数曲线（依坐标轴的选取为 \(y = a^{x}\) 或 \(y = \log x\)）的出现、对数螺线、双曲线的求积、\(\log(1 + x)\) 的级数展开，甚至记号 \(\log x\) 或 lx 的采用。至于三角函数，尽管在某种意义上它们可追溯到古代，值得注意的是，正弦曲线起初并不是作为由方程 \(y = \sin x\) 定义的曲线出现，而只是随罗贝瓦尔（{[}263{]}，第 63-65 页）作为“轮盘（roulette）的伴侣”出现（事实上，所涉及的正是曲线
\end{enumerate}

\[
y = R \left(1 - \cos \frac {x}{R}\right),
\]

也就是说，作为一条其定义由摆线的定义导出的辅助曲线。要遇到解析表达式的普遍概念，我们必须求诸 J. 格雷戈里，他在 1667 年（{[}169 a{]}，第 413 页）把它定义为这样一个量：它从其他的量出发，经由一系列代数运算“或任何别的可以想象的运算”而得到；他还在其序言中（{[}169 a{]}，第 408-409 页）试图精确化这一观念，办法是解释有必要在代数的五种运算\(^{14}\)之外再添上第六种运算，而这第六种运算最终无非就是取极限。但这些有趣的思想很快就被遗忘了，淹没在格雷戈里本人、牛顿以及其他人类发现的级数展开的洪流之中；这最后一种方法的巨大成功，在能有解析定义的函数与可展开为幂级数的函数之间造成了一种经久不消的混淆。

至于莱布尼茨，他似乎停留在笛卡尔的观点上，只是明确地附加了求积，又隐含地附加了他那个时代的分析中人们熟悉的其他运算：幂级数的求和、微分方程的求解。约翰·伯努利也是一样，当他想考虑 x 的任意函数时，把它引进为“一个以任意方式由 x 和常量构成的量”（{[}198 a{]}，第 III 卷，第 150 页），有时还明确说这是“以代数方式或超越方式”构成的一个量（{[}198 a{]}，第 III 卷，第 324 页）：而且在 1698 年，他与莱布尼茨商定给这样一个量以“x 的函数”之名（{[}198 a{]}，第 III 卷，第 507-510 和 525-526 页）。\(^{15}\)莱布尼茨此前已经引进了“常量”、“变量”、“参数”这些词，并在关于包络的讨论中明确了依赖于一个或多个参数的曲线族的概念（{[}198 a{]}，第 V 卷，第 266-269 页）。关于记号的问题在与约翰·伯努利的通信中也变得明确了：后者习惯用 X、或 \(\xi\) 来记 x 的任意函数（{[}198 a{]}，第 III 卷，第 531 页）；莱布尼茨表示赞同，但也提出用 \(\overline{x}|_{-1}^{1}\)、\(\overline{x}|_{-1}^{2}\)，而我们今天则会写作 \(f_{1}(x)\)、\(f_{2}(x)\)；他还为 x 的函数 z 的导数 \(\frac{dz}{dx}\) 提出记号 \(\bar{d}z\)（与表示微分的 \(dz\) 相对照），而伯努利则写作 \(\Delta z\)（{[}198 a{]}，第 III 卷，第 537 和 526 页）。

就这样，随着这个世纪的结束，英雄时代落幕了。新演算连同它的概念和记号，以莱布尼茨赋予它的形式确立起来。最早的弟子雅各布·伯努利和约翰·伯努利与大师竞相发现，自由地探索他为之指明道路的那片富饶疆域。第一部关于微分和积分演算的论著写于 1691 年和

1692 年，出自约翰·伯努利之手，\(^{16}\)供一位表现出好学生素质的侯爵之用。至于牛顿终于在 1693 年决定吝啬地发表对他的流数的一个简短一瞥（{[}326 a{]}，第 II 卷，第 391-396 页），这也无关紧要；即使他的《Principia》足以供一个多世纪的沉思，在无穷小演算的地盘上，他已被追上，并且在许多方面被超越。

此外，新体系的弱点至少在我们眼里是显而易见的。牛顿和莱布尼茨一举废除了延续两千年的传统，把首要角色给予了微分，并把积分贬为仅仅是它的逆运算；需要整个 19 世纪以及 20 世纪的一部分，才能通过把积分置于实变量函数的一般理论及其现代推广的基础之上，重新确立一种公正的平衡（见第 223 页）。也是从这种观点的颠倒，产生了不定积分以牺牲定积分为代价所扮演的过分、而且几乎是排他的角色——这一角色在巴罗那里已经出现，尤其是从牛顿和莱布尼茨开始：在这方面，同样是 19 世纪才把事情放回原位。最后，莱布尼茨特有的对符号作形式处理的倾向，在整个 18 世纪还将继续得到强调，远远超出当时分析的资源所能容许的程度。特别是必须承认，莱布尼茨的微分概念说实话没有任何意义；在 19 世纪初，它陷入了一种名誉扫地的境地，只是逐渐才从中恢复过来；而如果说一阶微分的使用最终得到了完全的正名，那么高阶微分尽管十分有用，直到我们今天也还没有真正恢复名誉。
